% Generated draft: requires mathematical and source review.
\documentclass[11pt]{article}
\usepackage[margin=25mm]{geometry}
\usepackage{fontspec}
\setmainfont{Latin Modern Roman}
\usepackage{amsmath,amssymb,mathtools,unicode-math}
\setmathfont{Latin Modern Math}
\usepackage{enumitem,needspace,xurl,hyperref}
\setlength{\parindent}{0pt}
\setlength{\parskip}{5pt}
\setlength{\emergencystretch}{3em}
\newcommand{\N}{\mathbb{N}}
\newcommand{\Z}{\mathbb{Z}}
\newcommand{\Q}{\mathbb{Q}}
\newcommand{\R}{\mathbb{R}}
\newcommand{\C}{\mathbb{C}}
\newcommand{\F}{\mathbb{F}}
\newcommand{\A}{\mathbb{A}}
\newcommand{\PP}{\mathbb{P}}
\newcommand{\E}{\mathbb{E}}
\newcommand{\eps}{\varepsilon}
\newcommand{\dd}{\,\mathrm{d}}
\begin{document}
\section*{UnsolvedMath ID: 30000964}
\textbf{Source problem code:} catalogue entry
\subsection{Definitions and mathematical statement}
Let \(Z\) be the counting measure of the chaotic analytic zero process from the report, whose points are zeros of a Gaussian random entire function. For an admissible test function \(h_r\), the customary centered linear statistic is
\[
Z_0(h_r)=\int h_r\,\dd Z-\E\int h_r\,\dd Z.
\]
The macroscopic convention often has \(h_r(z)=h(z/r)\); the record does not supply this definition, the precise Gaussian entire-function law, or the admissible class of \(h\), so these remain to be confirmed rather than imposed as source facts.

The corrected question asks whether the Gaussian moderate-deviation logarithmic tail asymptotic stated in the report continues to hold in the intermediate regime
\[
\frac{r}{(\log r)^2}\ll c\ll r\qquad(r\to\infty).
\]
With the intended asymptotic-separation reading of \(\ll\), this concerns choices \(c=c(r)>0\) satisfying
\[
\frac{c(r)(\log r)^2}{r}\longrightarrow\infty,
\qquad \frac{c(r)}r\longrightarrow0.
\]
The displayed lower endpoint is \(r/(\log r)^2\), not \((\log r)^2\). The supplied source-verification note explicitly corrects the latter expression in the raw record.

The tail event, variance normalization, coefficient of the Gaussian exponent, test-function hypotheses, and pointwise versus uniform interpretation across \(c(r)\) are not reproduced. Consequently no specific tail formula is substituted for the report's referenced asymptotic.

The raw record also asks a broader moderate-deviation and associated linear-response question for stationary random fields under hypotheses such as exponential correlation decay and finite exponential moments. The verification note identifies that as a separate source item, represented locally by OWR-1970-004, and it is not merged into the main chaotic-zero question here.
\subsection{Short English statement}
Does the report's Gaussian logarithmic tail asymptotic for centered chaotic-zero statistics extend to \(r/(\log r)^2\ll c\ll r\)? The corrected lower scale is retained; the precise statistic and tail formula need recovery.
\subsection{Sources}
\begin{itemize}
\item Local UnsolvedMath registry, permanent ID \(30000964\), source code catalogue entry, ``Moderate Deviations for Chaotic Analytic Zeros.'' No individual proposer is supplied.
\item Adjacent local records OWR-1970-001 (permanent ID \(30000963\)) describe the Gaussian entire-function zero process, and OWR-1970-004 (permanent ID \(30000965\)) carries the separate stationary-field linear-response question. They do not supply this record's complete tail asymptotic.
\item \emph{Stochastic Analysis} (2009), Oberwolfach Reports, DOI 10.4171/owr/2008/25, pp. 1361--1424: \url{https://doi.org/10.4171/owr/2008/25}. This is the sole supplied external reference.
\item Supplied literature triage, checked 2026-08-21: status ``open.'' It reports a moderate-deviation and linear-response question under exponential correlation decay and finite exponential moments, later restricted-model results, and no located general theorem under the stated hypotheses. The verification note separates that broad question from the corrected chaotic-zero item. This dated assessment derives from supplied OWR curation; time-sensitive status and cited claims require independent verification.
\end{itemize}
\end{document}
